\documentclass{report}

\usepackage[margin=0.5in]{geometry}
\usepackage{amsmath, mathtools, array, booktabs, hyperref, amsfonts}

\title{Homework 2}

\author{Darin Peries}

\date{\today}

\begin{document}
\maketitle

\section*{Problem 1}

\subsection*{Part 1}

To prove that any $A\in\mathbb{C}^{n\times n}$ follows the relation below, 
$$n^{-1} \| A \|_2 \leq n^{-1/2} \| A \|_{\infty} \leq \| A \|_2 \leq n^{1/2} \|A\|_1 \leq n \|A\|_2$$
In order to prove that the relation holds for all $A\in\mathbb{C}^{n\times n}$,
we need to rewrite the norm as $\|A\|=\sup\{\|Au\|, \ \|u\|=1\}=\frac{\|Au\|}{\|u\|}$.
The next step will involve the Holder inequality. This involves splitting the norm into two
norms being multiplied, resulting in, for example the 1-norm, 
$$\|Au\|_1 \leq  (\sum\limits_{i=1}^n |Au_{i}|^2)^{1/2} \times (\sum\limits_{i=1}^n |1|^2)^{1/2}=\sqrt{n}\|Au\|_2 $$ 
From the above expression, we can say that $\sup\{\|Au\|_1\} \leq \sqrt{n}\sup\{\|Au\|_2\}$ or $\|A\|_1 \leq \sqrt{n}\|A\|_2$.
This gives us a part of the relation. The idea can then be extended to both the 2-norm and infinity norm as follows. 
$\|Au\|^2_2=(\sqrt{\sum\limits_{i=1}^n|Au_i|^2})^2 \times 1 \leq \||Au|^2\|_{\infty} \|1\|_1=n\||Au|^2\|_{\infty}$. Which can then
be rewritten as $\|A\|_2 \leq \sqrt{n}\|A\|_{\infty}$ by taking the square root of both sides. For the 2-norm and 1-norm,
we know that the $\|A\|_{\infty} \leq n\|A\|_1$ and an upper bound of the 2-norm is $\sqrt{\|A\|_1\|A\|_{\infty}}$  
via the properties of the spectral radius covered in class, so we can then say $\|A\|_2 \leq \sqrt{\|A\|_1\|A\|_{\infty}}\leq \sqrt{\|A\|_1n\|A\|_1}=\sqrt{n}\|A\|_1$. 
Lastly, $\|A\|_{\infty} \leq n\|A\|_1 \leq n \sqrt{n} \|A\|_2$. This gives us the following
terms:
\begin{enumerate}
    \centering
    \item $\|A\|_1 \leq \sqrt{n} \|A\|_2$
    \item $\|A\|_2 \leq \sqrt{n} \|A\|_{\infty}$
    \item $\|A\|_2 \leq \sqrt{n} \|A\|_1$
    \item $\|A\|_{\infty} \leq n^{3/2} \|A\|_2$
\end{enumerate}

\noindent By scaling the bounds appropriately, say dividing by $n$, the relation can be upheld.

$$n^{-1}(\|A\|_2 \leq n^{1/2} \|A\|_{\infty}) \leq n^{-1/2}(n^{3/2}\|A\|_{2} \leq \sqrt{n} \|A\|_{1}) \leq n\|A\|_2$$ 

\subsection*{Part 2}

To prove that $\|x\| = \max\limits_{i=1,\dots,n}w_i|x_i|$ is a norm on $\mathbb{R}^n$, we need to show that it holds
the following:
\begin{enumerate}
    \item $\|x+y\| \leq \|x\| + \|y\|$
    \item $\|\lambda x\| = |\lambda|\|x\|$
    \item $\|x\| = 0 \iff x = 0$
\end{enumerate}

\subsubsection*{Triangle inequality}

$\|x+y\|=\max\limits_{i=1,\dots,n} w_i|x_i + y_i| = \max\limits_{i=1,\dots,n} |w_ix_i + w_iy_i| \leq \max\limits_{i=1,\dots,n} |w_ix_i| + \max\limits_{i=1,\dots,n} |w_iy_i| = \|x\| + \|y\|$
, via the subadditivity of the max function.

\subsubsection*{Positivity}

A norm on $\mathbb{R}^n$ needs to be $\geq 0$. In order to prove this, the first step involves assuming $x=0$. 
By assuming $x=0$ where $x\in\mathbb{R}^n$, the resulting norm is as follows:

$$\max\limits_{i=1,\dots,n} (w_i * 0, w_{i+1} * 0, \dots, w_n * 0) = 0$$

\noindent This shows that $\|x\|=0 \Longrightarrow x=0$. Now the second step assumes that $\|x\|=0$ and we need to prove that $x=0$.
Looking at the norm, $\max\limits_{i=1,\dots,n} w_i\|x_i\|$ can only be negative if $w_i$ is negative. Since $w_i$ is a positive
number, the norm being negative is not possible. Further more, the max function garuntees that the norm is non zero as long as some
$x_i>0$ for all $i$. Therefore, in order for $\|x\|=0$, $x=0$.

\subsubsection*{Homogeniety}

Given some scalar $\lambda\in\mathbb{R}$,
\[
\begin{aligned}
\|\lambda x\|
&= \max_{i=1,\dots,n} w_i|\lambda x_i|
= \max_{i=1,\dots,n} |\lambda|w_i|x_i| \\
&= |\lambda|\max_{i=1,\dots,n} w_i|x_i|
= |\lambda|\|x\|.
\end{aligned}
\]

\noindent Because the weights are all positive integers, and the max is taken over all the absolute values of x, the norm is greter than or equal to zero.
Since all three postiulates have been proven, $\|x\| = \max\limits_{i=1,\dots,n}w_i|x_i|$ is a norm on $\mathbb{R}^n$.

\subsubsection*{Subordinate Matrix Norm}
A subordinate matrix norm is the largest factor by which matrix can stretch a vector measured by a vetor norm.

\section*{Problem 2}

To prove that $\|A\|_1:=\max\limits_{x\neq 0} \frac{\|Ax\|}{\|x\|} = \max\limits_{j=1,\dots,n} \sum\limits_{i=1}^n |a_{ij}|$,
the definition of an induced matrix norm needs to be re-visited. 
$$\|A\| = \sup\limits_{\|u\|=1}\{Au\}$$ 

\noindent By definition of a supremum, the induced norm of $A$ needs to be greater than or equal to all elements in the resulting dot product for any unit vector $u$, ie an upper bound.
In order to find this, let $u$ be a vector of all zeros except for some index $j$ being equal to $1$ and thus satisfies $\|u\|_1=1$. By computing the dot product of $Au$, the $j^{th}$ column is 
extracted. The norm of the resulting vector is $\sum\limits_{i=1}^n |Au_i|$. In order to find an upper bound, multiply the matrix by a set of vectors $u$ where each vector in the set
is zero except some index $j$ equal to $1$. No vectors in the set have the same non zero $j$. For each vector, compute the dot product
and the 1-norm. Once all norms have been calculated, take the maximum norm as the upper bound of the induced matrix 1-norm.
Therefore, the upper bound, or supremum, of $\|A\|_1$ is $\max\limits_{j=1,\dots,n} \sum\limits_{i=1}^n |a_{ij}|$.

\section*{Problem 3}
\subsection*{Part 1}
To show that $A$ is non-singular by $\|Ax\| \geq \theta \|x\|, \ \forall x$, first set $x$ to zero. Since $\|.\|$ is a norm on $\mathbb{R}^n$,
$x$ needs to be zero in order for the norm to be zero. This means that the $kern(A)=\{0\}$. Because the the kernel of $A$
is only the zero vector, the matrix is invertible, a.k.a non-singular. \\

\noindent Since $A$ is invertible, to show that $\|A^{-1}\|\leq \theta^{-1}$, say there is some vectopr $y$, where $Ax=y$ and $A^{-1}y=x$.
Since $\|Ax\| \geq \theta\|x\|$, then $\|y\| \geq \theta \|A^{-1}y\|$. Rearranging both sides, the expression is as follows, 
$\frac{1}{\theta} \geq \frac{\|A^{-1}y\|}{\|y\|}$, which can be rewritten as $\frac{1}{\theta} \geq \sup\{\|A^{-1}y\|\}$, or more simply as
$\theta^{-1} \geq \|A^{-1}\|$.

\subsection*{Part 2}
To show that $A$ is invertible based on $\|I-P(A)\|\leq1$, the Neuman series can be used. This is because the norm of $A-I$ is less than $1$. The series is as follows:

$$(I-A)^{-1} = \sum\limits_{i=0}^{\infty} A^i$$

\noindent Substituting $A$ with $I-P(A)$ gives $P(A)^{-1}=\sum\limits_{i=0}^{\infty}(I-P(A))^i$. Since $P$ is a polynomial with no constant terms,
$A$ can be factored out of the matrix, resulting in $(AG(A))^{-1} = \sum\limits_{i=0}^{\infty} (I-P(A))^i$, where $G(A)$ is the resulting polynomial 
with one constant term after $A$ is factored out. Rewritting the expression, $A^{-1} = G(A)\sum\limits_{i=0}^{\infty}(I-P(A))^i$.

\section*{Problem 4}
\subsection*{part 1}
If $A$ is positive symmetric definite, then its diagonals must be $>0$. Say there is a unit vector $e_k$, where all entries are zero except for the entry at 
index k, which is set to $1$. using the quadratic form, $<e_kA,\ e_k^T>$, The only influencing entry in the matrix that determines whether the scalar result is positive is
the $k^{th}$ diagonal. If the quadratic form is computed from $k=1$ to $n$, the products should all be $>0$.

\subsection*{part 2}
Say there is some vector $x\in\mathbb{R}^n$ and on index $i$, $x_i=\alpha$ and on index $j$, $x_j=1$ with zeroes everywhere else.
By using the quadratic form $<Ax, \ x^T>$, the following vector and equation are generated.

$$Ax = ((\alpha a_{1i} + a_{1j}), \dots ,(\alpha a_{ii} + a_ij), \dots, (\alpha a_{ji} + a_jj)), \dots, (\alpha a_{ni}+a_{nj})$$

$$<Ax, x^T> = \alpha^2 a_{ii} + \alpha a_{ij} + \alpha a_{ji} + a_{jj} = \alpha^2 a_{ii} + 2\alpha a_{ij} + a_{jj}$$

\noindent When factoring the equation above, the roots of the polynomial are $a_{ii}(\alpha+\frac{a_{ij}}{a_{ii}})^2+(a_{jj}+\frac{a_{ij}^2}{a_{jj}})=(a_{jj}-\frac{a_{ij}^2}{a_{ii}}) > 0$.
Rearranging the equation gives $a_{ii}a_{jj} > a_{ij}^2$.
\end{document}